\clearpage\begingroup\SourceLeaf{405b}\ArithmeticHead{74}{L}{STÄCKEL, GAUSS ALS GEOMETER.}
\renewcommand{\thefootnote}{\arabic{footnote})}\setcounter{footnote}{1}
\begin{center}
\textbf{24.}\par\smallskip
\textls[100]{Die Vielecke.}
\end{center}
\noindent\hspace*{\parindent}Durch eine Anfrage von \textsc{Schumacher} vom 19. März 1836 veranlasst
(W. X\textsubscript{1}, S. 459) hat sich \textsc{Gauss} mit der Frage nach dem »kürzesten Verbindungssystem«
von beliebig vielen, im Besonderen von vier Punkten beschäftigt
(W. X\textsubscript{1}, S. 461—467), einer glücklichen Verallgemeinerung der
Summe der Entfernungen eines Punktes von gewissen gegebenen Punkten,
die noch heute eingehendere Erforschung verdiente\footnote{Vgl. auch die Dissertation von K. \textsc{Bopp}, \textit{Das kürzeste Verbindungssystem von vier Punkten}, Göttingen 1879.}.
\par
Die Lösung der Aufgabe, in einen gegebenen Kreis ein Vieleck zu beschreiben,
dessen Seiten durch je einen gegebenen Punkt gehen, ist wieder
der Verwendung komplexer Grössen entsprungen (W. IV, S. 398, Zusatz V,
S. 369).
\par
Im ersten Abschnitt (Nr. 10) ist bemerkt worden, dass \textsc{Gauss} die Frage
aufgeworfen hat, was man unter dem Inhalt eines beliebigen Vielecks zu
\clearpage\SourceLeaf{406}\ArithmeticHead{75}{R}{ELEMENTARE UND ANALYTISCHE GEOMETRIE.}\setcounter{footnote}{0}
\noindent verstehen habe. Auf den Inhalt eines Vielecks bezieht sich die folgende Stelle
in dem Zusatz I zu \textsc{Schumachers} Übersetzung der \textit{Géométrie de position} von
\textsc{Carnot}\footnote{Werke XI\textsubscript{1}.} (S. 362):
\par
\textls[100]{Anmerkung des Herausgebers} [\textsc{Schumacher}]. »Es ist, nach einem
schönen Theorem des Herrn Professor \textsc{Gauss}, der Inhalt eines Vielecks von
\(n\) Seiten, wenn die Koordinaten der Winkelpunkte nach der Reihe in einer
Richtung gezählt:
\[
x,y;\ x',y';\ \ldotp\,\ldotp\,\ldotp\ x^{(n-1)},y^{(n-1)}
\]
sind,
\[
=\tfrac12\{x(y'-y^{(n-1)})+x'(y''-y)+x''(y'''-y')+\ldotp\,\ldotp\,\ldotp+x^{(n-1)}(y-y^{(n-2)})\},
\]
worüber Er selbst vielleicht, bey einer andern Gelegenheit, uns eine vollständige
Abhandlung schenken wird«.
\par
Die Ankündigung einer vollständigen Abhandlung macht es wahrscheinlich,
dass \textsc{Gauss} schon damals die Verallgemeinerung auf beliebige Vielecke
im Auge hatte, bei denen also der Umfang sich selbst durchsetzen kann, wie
er sie in dem Brief an \textsc{Olbers} vom 30. Oktober 1825 andeutet (W. VIII,
S. 398)\footnote{In dem Brief an \textsc{Olbers} vom 30. Oktober 1825 bemerkt \textsc{Gauss}, er habe »erst vor kurzem eine Abhandlung von \textsc{Meister} im ersten Bande der Novi Commentarii Gotting. kennen gelernt, worin die Sache fast ganz auf gleiche Art betrachtet und sehr schön entwickelt wird«; gemeint ist die Abhandlung von A. L. F. \textsc{Meister}, \textit{Generalia de genesi figurarum planarum et inde pendentibus earum affectionibus}, Novi Commentarii acad. Gotting., vol. I ad annos 1769/70, 1771, S. 144.}. In einer aus dem Nachlass 1866 herausgegebenen Abhandlung hat
\textsc{Jacobi} eine Regel für die Bestimmung des Inhalts gegeben\footnote{C. G. J. \textsc{Jacobi}, \textit{Regel zur Bestimmung des Inhalts der Sternpolygone}, Journal für die r. u. a. Mathematik Bd. 65, 1866, S. 173, Werke VII, S. 40; vgl. auch W. \textsc{Veltmann}, \textit{Berechnung des Inhalts eines Vielecks aus den Coordinaten der Eckpunkte}, Zeitschrift für Mathematik und Physik, Bd. 32, 1887, S. 339. Nach L. \textsc{Königsberger}, \textit{C. G. J. Jacobi}, Leipzig 1904, S. 155 hat \textsc{Jacobi} seine Regel im Sommer 1833 gefunden.}.
\par
In der Behaftung von Inhalten mit Vorzeichen ist \textsc{Möbius} mit \textsc{Gauss}
zusammengetroffen, zuerst im \textit{Barycentrischen Calcul} (1827)\footnote{A. F. \textsc{Möbius}, \textit{Der barycentrische Calcul}, Leipzig 1827, Kap. II, § 17 und 18; Werke I, S. 39—41.}, dann in der Abhandlung
\textit{über den Inhalt der Polyeder} (1865)\footnote{A. F. \textsc{Möbius}, \textit{Über die Bestimmung des Inhaltes eines Polyeders}, Leipziger Berichte, Bd. 17, 1865, S. 31, Werke II, S. 485—491.}. Wenn die Mathematiker des
20. Jahrhunderts diese Dinge als selbstverständlich ansehen, so hat doch
\par\smallskip\noindent\hfill{\small 10\textsuperscript{*}}\par
\clearpage\SourceLeaf{407a}\ArithmeticHead{76}{L}{STÄCKEL, GAUSS ALS GEOMETER.}\setcounter{footnote}{0}
\noindent \textsc{Baltzer}, in den Erinnerungen an die \textsc{Gauss}schen Zeiten wurzelnd, mit Recht
hervorgehoben, dass die Bestimmung des Zeichens einer Strecke nach der
voraus bestimmten positiven Richtung einer Geraden, einer Dreiecksfläche
nach dem voraus bestimmten positiven Sinn ihrer Ebene und eines Tetraederinhalts
nach einem voraus bestimmten Schraubungssinn beim Erscheinen des
\textit{Barycentrischen Calculs} »neu und fast befremdend« erschienen seien\footnote{R. \textsc{Baltzer}, \textit{Vorrede über Möbius}, \textsc{Möbius}’ Werke I, S. VIII.}.
\par
\textsc{Gauss} hat auch eine von \textsc{Möbius} gestellte Aufgabe\footnote{A, F. \textsc{Möbius}, \textit{Beobachtungen auf der Sternwarte zu Leipzig usw.}, Leipzig 1823, S. 57; Werke I, S. 394.} gelöst, die besagt,
man solle den Inhalt eines Fünfecks aus den Inhalten der fünf Dreiecke bestimmen,
die von den Verbindungsstrecken der fünf Eckpunkte gebildet werden
(W. IV, S. 406).
\clearpage\endgroup
